\section{Experimental Evaluation}
\label{sec:experimental_evaluation}

We evaluate the implemented test-generation workflow using policy
fragments from three application domains: healthcare access control,
financial management, and student information management. The
experiments examine test generation, agreement between two executable
representations, and detection of injected policy faults. The current
prototype reads concrete bindings exported as Event-B contexts; it does
not execute the original Event-B machine events. Accordingly, the
results constitute a concrete-data feasibility evaluation rather than
an end-to-end evaluation from verified models to an independent PDP.

\subsection{Experimental Setup}
\label{subsec:exp_setup}

\paragraph{Prototype and environment.}
The prototype is implemented in Python~3.12.14 and runs on Linux
x86\_64 with kernel~6.18.44 and glibc~2.39. It uses the Python standard
library, including \texttt{xml.etree.ElementTree} for XML processing.
The tool parses the supplied rule fragments, constructs concrete
Event-B bindings, derives finite request domains, selects test requests,
and exports XACML policy/request XML and expected decisions. Results
are recorded as CSV and JSON. Experiments are sequential, with fixed
attribute ordering and deterministic witness selection. No stochastic
optimization or post-generation suite reduction is applied to the
three datasets.

Expected decisions are computed by a reference evaluator using bindings
extracted from the generated context. A separately implemented local
XML-subset evaluator processes the emitted policies. For HACS and SIS,
this evaluator compiles equality targets for exhaustive evaluation;
selected serialized requests are additionally evaluated through an
uncompiled XML evaluator. These evaluators share the study's semantic
assumptions. No independent production PDP, XACML schema validation,
Rodin proof replay, or ProB execution is included in these measurements.

\paragraph{Healthcare Access Control System (HACS).}
HACS represents authorization requests concerning healthcare-domain
resources. The supplied dataset contains 50 rule fragments, with 25
Permit effects and 25 Deny effects. Targets specify a subject, a
resource, and an action, represented by values such as
\texttt{Admin\_1}, \texttt{File\_01}, and \texttt{read}.
The fragments use a simplified \texttt{PolicyRule} representation;
they are normalized into XACML~3.0 rules whose targets are conjunctions
of singleton string-equality matches. Policy membership from the
repository organizes these rules into 20 policies. Because the supplied
fragments do not declare a combining algorithm, DenyOverrides is used
as an explicit experimental assumption.

\paragraph{Financial Management System (FMS).}
FMS represents authorization requests in a financial-management
setting. The supplied dataset contains 45 Permit rules organized into
18 policies using DenyOverrides. Fifteen policies contain three rules
each, while the remaining three policies are empty. A textual target
such as \texttt{Role1 / Action1} is interpreted as the conjunction
\texttt{role = Role1} and \texttt{action = Action1}. Empty conditions
are interpreted as true. The empty policies are retained and evaluated
as NotApplicable. The generic role/action names provide a structural
benchmark rather than evidence of a deployed financial policy.

\paragraph{Student Information System (SIS).}
SIS represents access to student-information resources. Its 40 rule
fragments contain 20 Permit and 20 Deny effects, with subjects such as
\texttt{Teacher\_1}, resources such as \texttt{Grades\_01}, and actions
\texttt{view} and \texttt{modify}. As in HACS, the simplified fragments
are normalized into equality-based XACML targets. Repository membership
organizes the rules into 15 policies. DenyOverrides is assumed because
the supplied fragments do not specify the combining algorithm.

\paragraph{Data provenance and interpretation.}
The three rule datasets are supplied as JSON. HACS and SIS membership
metadata is obtained from the XACML2EventB repository at commit
\texttt{8cc1782f3df3}.\footnote{\url{https://github.com/binhtt/XACML2EventB}}
Placeholder XML policy strings in the repository are not executed.
All policy targets are assumed universal; omitted conditions are true.
Environment values are singleton and do not affect applicability.
FMS additionally uses an irrelevant singleton resource domain.
Each policy is tested separately; policy-set aggregation is not
included. Table~\ref{tab:exp_datasets} summarizes the resulting instances.

\begin{table}[htbp]
\centering
\caption{Characteristics of the concretized datasets}
\label{tab:exp_datasets}
\small
\begin{tabular}{lrrrrr}
\hline
System & Rules & Policies & Permit & Deny & Empty policies \\
\hline
HACS & 50 & 20 & 25 & 25 & 0 \\
FMS & 45 & 18 & 45 & 0 & 3 \\
SIS & 40 & 15 & 20 & 20 & 0 \\
\hline
Total & 135 & 53 & 90 & 45 & 3 \\
\hline
\end{tabular}
\end{table}

\subsection{Test Configuration and Evaluation Criteria}
\label{subsec:exp_configuration}

\paragraph{Concrete bindings and request domains.}
The prototype exports concrete definitions of rule targets, conditions,
effects, policy membership, and rule ordering into an Event-B context.
A restricted parser reads these generated formulas back and checks
correspondence with the input data. This parser supports the emitted
binding subset, not arbitrary Event-B expressions. The new contexts
have not been type-checked or proved in Rodin, and original machine
guards and actions are not interpreted by the current generator.

Each attribute domain contains the values appearing in the dataset
and one additional unmatched value. FMS uses 46 role values and 46
action values, producing $46\times46=2,116$ requests per policy.
HACS uses 51 subject values, 51 resource values, and five action values,
producing $51\times51\times5=13,005$ requests. SIS uses 41 subject
values, 41 resource values, and three action values, producing
$41\times41\times3=5,043$ requests. Requests contain all relevant
attributes, each with a single string value.

\paragraph{Test objectives and selection.}
The generator constructs objectives for applicable and non-applicable
outcomes of each rule and for feasible Permit, Deny, and NotApplicable
policy decisions. Complete enumeration identifies feasible objectives
within the configured domain. Objectives are processed in a fixed
order. If the selected suite does not already satisfy an objective,
the generator adds the first satisfying request in the configured
ordering. The resulting test contains a policy identifier, a request,
and its expected decision. Coverage is computed over feasible
rule-applicability and decision objectives; it is not Event-B event
coverage, MC/DC, or coverage of combining interactions.

\paragraph{Mutation protocol.}
Mutation analysis applies single-site edits to the emitted policy XML
while preserving the reference bindings and oracle decisions. All
three datasets use effect replacement, rule removal, and subject/role
value replacement. HACS and SIS additionally use resource-value and
action-value replacements. Value replacements use the corresponding
unmatched domain value. Effect replacement exchanges Permit and Deny.
Two combining mutants per policy replace DenyOverrides with
PermitOverrides and FirstApplicable, respectively.

A mutant is detected when a selected test produces a decision different
from the unchanged reference decision. Each mutant is also evaluated
on the complete finite request domain to identify domain-equivalent
variants. Mutation counts refer to generated operator instances;
behaviorally duplicate instances are not deduplicated.

\paragraph{Evaluation criteria.}
We report policy-specific test evaluations, distinct request values,
coverage, internal decision agreement, and mutation detection.
A request reused against different policies counts once per policy
in test evaluations. The raw mutation score includes all generated
mutants. The adjusted score excludes only mutants shown equivalent
on the declared finite domain; no global equivalence claim is made.

\subsection{Test-Generation and Decision-Agreement Results}
\label{subsec:exp_generation_results}

Table~\ref{tab:exp_tests} summarizes generation and reference checks.
The prototype selects 70 HACS, 63 FMS, and 55 SIS tests, for a total
of 188 policy-specific evaluations. These suites contain 51, 46,
and 41 distinct request values, respectively. All feasible
rule-applicability and policy-decision objectives are covered.

\begin{table}[htbp]
\centering
\caption{Test generation and internal reference agreement}
\label{tab:exp_tests}
\small
\begin{tabular}{lrrrr}
\hline
System & Tests & Distinct requests & Reference checks & Mismatches \\
\hline
HACS & 70 & 51 & 260,100 & 0 \\
FMS & 63 & 46 & 38,088 & 0 \\
SIS & 55 & 41 & 75,645 & 0 \\
\hline
Total & 188 & -- & 373,833 & 0 \\
\hline
\end{tabular}
\end{table}

Each nonempty policy requires one positive request per rule and an
additional request for the non-applicable case. The three empty FMS
policies each require one request. The extracted-binding and XML-subset
evaluators agree on all 373,833 policy/request combinations in the
configured domains. This agreement supports consistency of the
implemented representations under the recorded assumptions; it does
not establish correctness against an independent PDP or the complete
XACML specification.

\subsection{Fault-Detection Results}
\label{subsec:exp_fault_detection}

Table~\ref{tab:exp_mutants} reports mutation outcomes. The suites
detect 250 of 290 HACS mutants, 135 of 171 FMS mutants, and 200 of
230 SIS mutants. Their raw mutation scores are 86.21\%, 78.95\%,
and 86.96\%, respectively. The pooled score is $585/691=84.66\%$.

\begin{table}[htbp]
\centering
\caption{Mutation outcomes on the finite request domains}
\label{tab:exp_mutants}
\small
\begin{tabular}{lrrrr}
\hline
System & Mutants & Detected & Domain-equivalent & Raw score \\
\hline
HACS & 290 & 250 & 40 & 86.21\% \\
FMS & 171 & 135 & 36 & 78.95\% \\
SIS & 230 & 200 & 30 & 86.96\% \\
\hline
Total & 691 & 585 & 106 & 84.66\% \\
\hline
\end{tabular}
\end{table}

All evaluated effect, rule-removal, and attribute-value mutants are
detected (Table~\ref{tab:exp_operators}). Positive requests provide
direct witnesses: changing an applicable rule's effect changes its
expected decision, while removing the rule or changing one of its
target values makes the original request non-applicable.

\begin{table}[htbp]
\centering
\caption{Detected/generated mutants by operator family}
\label{tab:exp_operators}
\small
\begin{tabular}{lrrr}
\hline
Operator & HACS & FMS & SIS \\
\hline
Effect replacement & 50/50 & 45/45 & 40/40 \\
Rule removal & 50/50 & 45/45 & 40/40 \\
Subject/role replacement & 50/50 & 45/45 & 40/40 \\
Resource replacement & 50/50 & -- & 40/40 \\
Action replacement & 50/50 & -- & 40/40 \\
Combining-algorithm replacement & 0/40 & 0/36 & 0/30 \\
\hline
\end{tabular}
\end{table}

For example, HACS\_POLICY\_1 contains a Permit rule matching
$(\texttt{Admin\_1},\texttt{File\_01},\texttt{read})$.
This request is included in the generated suite. Replacing the rule's
effect with Deny changes the decision to Deny, whereas removing the
rule changes it to NotApplicable. Both outcomes differ from the
reference Permit decision.

None of the 106 combining-algorithm mutants is detected. Exhaustive
checks establish their equivalence within the configured domains.
The targets in each dataset are mutually exclusive, so no request
activates multiple rules of a policy. In addition, FMS contains only
Permit effects. Although HACS and SIS include both effects, those
effects do not interact in a single evaluation. Excluding the
domain-equivalent mutants yields an adjusted score of 100\% for
each dataset. This score concerns the evaluated rule-level mutations
and should not be interpreted as general combining-fault detection.

\subsection{Discussion and Threats to Validity}
\label{subsec:exp_threats}

\paragraph{Model and oracle validity.}
Concrete bindings are generated and read back, but the original
Event-B event semantics are not executed. The new contexts have not
been proved in Rodin. The uploaded historical project also contains
confidence-zero and broken proof-status entries; its proofs were not
replayed. Agreement between local evaluators may conceal shared
interpretation errors. Schema validation and an independent PDP are
required before claiming standard-compliant executable requests or
end-to-end testing from verified specifications.

\paragraph{Dataset validity.}
The rule fragments are simplified, and HACS/SIS combining algorithms
and universal policy targets are assumed. Generic attribute names,
mutually exclusive targets, missing nontrivial conditions, and empty
FMS policies limit representativeness. No overlapping Permit/Deny
rules, Indeterminate outcomes, missing attributes, multivalued bags,
or root-policy-set aggregation are evaluated.

\paragraph{Mutation and comparative validity.}
The injected mutants are controlled operator instances rather than
naturally occurring defects. Domain equivalence does not imply
unrestricted semantic equivalence. The datasets make several mutations
straightforward to detect, and no alternative generator comparison,
ablation, or repeated timing study is included in these three-system
measurements. The results support feasibility and observed detection,
not comparative superiority or general correctness.

\paragraph{Reproducibility.}
The artifact includes input fragments, membership metadata, file hashes,
concrete context candidates, parser and evaluator sources, generated
policy/request XML, expected decisions, and per-policy CSV results.
Fixed domain and objective ordering support deterministic reproduction.
Historical Rodin proof-status records are kept separate from the
new test-generation and mutation measurements.
